\chapter{Teacher and Hierarchy Control Details}
\label{app:teacher-hierarchy-details}
\label{app:teacher-transfer-probe}

This appendix records the compact details behind the teacher-derived hierarchy
and reverse-order control summarized in Chapter~\ref{chap:vision-curriculum}.

\section{Transfer-teacher probe}

The teacher probe used CIFAR-100 with a ResNet-18 initialized from ImageNet
weights and fine-tuned for 30 epochs using SGD, learning rate 0.01, weight decay
\(5\times 10^{-4}\), and data augmentation. The best test accuracy was in the
77.65\%--78.14\% range across the two relevant probes, substantially higher
than the earlier from-scratch ResNet-18 baseline used in the model-size study.
This is a test-accuracy summary. The word ``validation'' below refers to the
held-out images used for embedding extraction, not to a claim that validation
accuracy was higher or lower than test accuracy. The teacher's validation
embeddings were used as the distance source for a learned class hierarchy.

The teacher hierarchy is not interpreted as a human ontology. It is a
representation-induced class partition. Some groups that look semantically loose
to a human can still be useful if they reflect visual statistics in CIFAR-100,
such as shape, texture, background, object scale, or typical viewpoint.

\section{Reverse-order hierarchy control}

The reverse-order condition in the main text is a hierarchy-order control. It
uses the same teacher-derived hierarchy and the same marginalized coarse-label
loss as the regular curriculum, but it visits the non-singleton hierarchy levels
in reverse order before switching to fine labels.

For a typical teacher hierarchy, the order is approximately:

\begin{center}
\begin{tabular}{lll}
\toprule
Regime & Stage order before fine-tuning & Meaning \\
\midrule
Easy-to-hard & 3--5 clusters $\rightarrow$ 23--24 clusters & coarse to less coarse \\
Reverse-order & 23--24 clusters $\rightarrow$ 3--5 clusters & less coarse to coarse \\
\bottomrule
\end{tabular}
\end{center}

Both regimes then finish with ordinary 100-class fine-label training. Therefore
the reverse-order condition is not a wrong-label control and not a pure negative
control. Its early performance can be relatively high because 23--24 clusters on
CIFAR-100 correspond to small class bundles, which are much closer to the final
100-way task than a 3--5 cluster level.

\section{Per-seed hierarchy-control results}

\begin{table}[h]
\centering
\small
\begin{tabular}{rrrrrr}
\toprule
Seed & Baseline & Random mean & Reverse-order & Self hierarchy & Teacher hierarchy \\
\midrule
42 & 35.53 & 36.41 & 38.10 & 37.15 & 37.48 \\
43 & 35.37 & 36.47 & 37.21 & 38.55 & 38.06 \\
44 & 35.23 & 36.16 & 36.99 & 37.59 & 37.52 \\
\midrule
Mean & 35.38 & 36.35 & 37.43 & 37.76 & 37.69 \\
\bottomrule
\end{tabular}
\caption{Per-seed best test accuracy for the hierarchy-control suite. Values are
percentages. Random mean averages the random hierarchy seeds within each
training seed.}
\label{tab:app-teacher-control-seeds-compact}
\end{table}

The per-seed table supports the main interpretation: random coarsening is a real
positive control, but structured hierarchies remain stronger on average. The
teacher hierarchy is competitive with the self-derived hierarchy but does not
clearly dominate it in peak accuracy.
